\section{Tension d'une pile médicale}
\label{ex:34-tension}
Dans le domaine médical, la pile $\ominus ~ \ce{Li} | \text{electrolyte} | \ce{I_2} ~ \oplus$ est utilisée comme source d'énergie pour les pacemakers. 

\Q Donner les demi-équations rédox associées à l'anode et la cathode, et en déduire l'équation de réaction.
\Q Calculer l'enthalpie standard de réaction à \SI{298}{\kelvin}.
\Q Calculer l'entropie standard de réaction à \SI{298}{\kelvin}.
\Q En déduire l'enthalpie libre de réaction à \SI{298}{\kelvin} puis déterminer la f.e.m dans ces conditions.

Dans ce dispositif, la tension de sortie a été mesurée en fonction de la quantité de charge déchargée. 

\begin{minipage}{0.55\textwidth} 
\includegraphics[width=\textwidth]{figures/Fig01}
\end{minipage} \hfill \begin{minipage}{0.40\textwidth} 
\captionof{figure}{\'{E}volution de la tension de sortie et de la résistance interne d'une batterie \ce{Li}/\ce{I2} utilisée dans les pacemakers.}
\label{fig:Chp1-01}
\end{minipage}

\Q Montrer l'accord entre vos grandeurs et la mesure sur la majorité du cycle de décharge de la pile. 
\Q En vous appuyant sur l'équation de réaction, lier l'écroulement du potentiel de réaction au-delà de \SI{1.2}{\ampere\hour} à la structure de la pile.

On se placera dans l'approximation d'Ellingham.
\Q En déduire la relation entre la variation de la force électromotrice $e^\circ$ avec la température.
\Q Commenter et justifier la valeur obtenue.

\textbf{Données}
\begin{center}
\begin{tabular}{ccccc}
\toprule
 & & Li & \ce{I2} & \ce{LiI} \\
\midrule
Enthalpie de formation & [\si{\kilo\joule\per\mole}] & 0 & 0 & -270.1\\
Entropie molaire & [\si{\joule\per\mole\per\kelvin}] & 29.08 & 116.14 & 85.77 \\
\bottomrule
\end{tabular}
%Table 2.1, advanced battery.
\end{center}